\documentclass[12pt]{article}
\usepackage[utf8]{inputenc}
\usepackage[T5]{fontenc}
\usepackage[vietnamese]{babel}
\usepackage[a4paper,margin=2.5cm]{geometry}
\usepackage{setspace}
\usepackage{indentfirst}
\usepackage{titlesec}
\usepackage{hyperref}
\usepackage{graphicx}
\usepackage{amsmath,amssymb,amsthm,mathtools}
\usepackage{enumitem}

\setstretch{1.25}
\setlength{\parindent}{1cm}
\setlength{\parskip}{0.45em}

\titleformat{\section}
  {\normalfont\Large\bfseries}{\thesection}{1em}{}
\titlespacing{\section}{0pt}{1.2em}{0.7em}

\hypersetup{
    colorlinks=true,
    linkcolor=blue,
    filecolor=magenta,
    urlcolor=cyan,
    citecolor=blue
}

\title{Mô hình hóa bài toán motion planning \\ thành bài toán optimal control trên đồ thị các vùng lồi bằng phương pháp multiple shooting}
\author{}
\date{}

\begin{document}

\maketitle
\tableofcontents
\newpage

\section{Mục tiêu}

Mục tiêu là formulate một bài toán motion planning \textbf{tìm đồng thời}
\begin{enumerate}[label=(\roman*)]
    \item một \textbf{chuỗi rời rạc các vùng an toàn lồi}, và
    \item một \textbf{quỹ đạo liên tục theo thời gian} thỏa động lực học,
\end{enumerate}
sao cho toàn bộ quỹ đạo tránh vật cản và tối ưu một hàm chi phí toàn cục.

Phần \emph{discrete} (chọn vùng nào, theo thứ tự nào) và phần \emph{continuous} (quỹ đạo, điều khiển, thời lượng) phải được tối ưu \textbf{cùng một lúc}.

Về nguồn gốc ý tưởng, mô hình này ghép hai khối lý thuyết:
\begin{itemize}
    \item \textbf{Graphs of Convex Sets / network flow} để chọn một đường đi rời rạc qua các vùng an toàn lồi \cite{marcucci2021,marcucci2022};
    \item \textbf{direct multiple shooting} để mô hình hóa quỹ đạo liên tục trong mỗi vùng lồi.
\end{itemize}

Trong \emph{Shortest Paths in Graphs of Convex Sets}, quyết định tối ưu bao gồm đồng thời một đường đi rời rạc trên đồ thị và các biến liên tục gắn với những đỉnh/cạnh được kích hoạt \cite{marcucci2021}. Trong \emph{Motion Planning around Obstacles with Convex Optimization}, không gian collision-free được phân hoạch thành các safe convex regions, và mục tiêu là xây dựng một continuous trajectory nằm hoàn toàn trong hợp của các vùng ấy \cite{marcucci2022}. Còn trong \emph{direct multiple shooting}, state trên mỗi đoạn thời gian được thay bởi nghiệm của một local initial value problem bắt đầu từ một initial value, và các đoạn được nối lại bằng continuity defects \cite{bockplitt,diehl2005}.

\section{Bài toán gốc cần giải}

Xét một hệ động lực học liên tục
\begin{equation}
\dot x(t)=f(x(t),u(t)),
\qquad t\in[0,T],
\label{eq:ct_dyn}
\end{equation}
trong đó:
\begin{itemize}
    \item $x(t)\in\mathbb R^{n_x}$ là trạng thái,
    \item $u(t)\in\mathbb R^{n_u}$ là điều khiển,
    \item $T>0$ là thời điểm kết thúc (có thể cố định hoặc là biến quyết định).
\end{itemize}

Cho trạng thái đầu và trạng thái đích
\[
x(0)=x^{\mathrm s},
\qquad
x(T)=x^{\mathrm g},
\]
và cho không gian tự do (collision-free space)
\[
Q \subset \mathbb R^{n_x}.
\]

Trong toàn bộ báo cáo, để tránh nặng ký hiệu, ta viết trực tiếp $x(t)\in Q$. Nếu trạng thái $x$ còn chứa thêm các biến động lực không mang ý nghĩa hình học, thì điều kiện collision-free đúng cần được hiểu là $\pi(x(t))\in Q$ đối với một phép chiếu cấu hình $\pi$; các vùng $Q_v$ về sau cũng được hiểu theo cùng nghĩa đó.

Ta muốn giải bài toán optimal control tổng quát
\begin{align}
\min_{T,\,x(\cdot),\,u(\cdot)}\quad
& \Phi\bigl(x(T),T\bigr)+\int_0^T \ell\bigl(x(t),u(t)\bigr)\,dt
\label{eq:orig_ocp_obj}\\
\text{s.t.}\quad
& \dot x(t)=f(x(t),u(t)), && t\in[0,T],
\label{eq:orig_ocp_dyn}\\
& x(0)=x^{\mathrm s},\qquad x(T)=x^{\mathrm g},
\label{eq:orig_ocp_bc}\\
& x(t)\in Q,\qquad u(t)\in \mathcal U, && \forall t\in[0,T].
\label{eq:orig_ocp_path}
\end{align}


Trong hàm mục tiêu \eqref{eq:orig_ocp_obj}, 
$\Phi(x(T),T)$ là chi phí cuối (terminal cost), dùng để đánh giá trạng thái đích và/hoặc thời gian kết thúc; 
còn $\ell(x(t),u(t))$ là chi phí tức thời (running cost), được tích lũy trên toàn khoảng thời gian thông qua tích phân 
$\int_0^T \ell(x(t),u(t))\,dt$. 
Vì vậy, hàm mục tiêu gồm một phần đánh giá tại điểm cuối và một phần đánh giá toàn bộ quá trình chuyển động.

\section{Phân hoạch không gian tự do thành các vùng an toàn lồi}

Giả sử ta đã có một họ hữu hạn các vùng an toàn lồi
\[
\mathcal Q:=\{Q_v\}_{v\in V_r},
\qquad
Q_v\subseteq Q \subset \mathbb R^{n_x},
\]
trong đó mỗi $Q_v$ là một tập lồi đóng, bị chặn, không giao với vật cản.

Từ các vùng này, ta xây dựng một đồ thị có hướng
\[
G=(V,E),
\qquad
V=V_r\cup\{\sigma,\tau\},
\qquad
E=E_{rr}\cup E_s\cup E_t,
\]
trong đó:
\begin{itemize}
    \item $V_r$ là tập các đỉnh tương ứng với vùng lồi,
    \item $\sigma$ là đỉnh nguồn,
    \item $\tau$ là đỉnh đích.
\end{itemize}

Các cạnh giữa hai vùng lồi được tạo theo điều kiện adjacency hình học
\begin{equation}
E_{rr}:=\{(u,v)\in V_r\times V_r:\ Q_u\cap Q_v\neq \varnothing\}.
\label{eq:adjacency}
\end{equation}

Các cạnh biên được thêm riêng cho source và target:
\begin{equation}
E_s:=\{(\sigma,v):\ x^{\mathrm s}\in Q_v\},
\qquad
E_t:=\{(v,\tau):\ x^{\mathrm g}\in Q_v\}.
\label{eq:boundary_edges}
\end{equation}

Điều kiện \eqref{eq:adjacency} chỉ nói rằng hai vùng chạm nhau về mặt hình học, tức là tồn tại một trạng thái kết nối thuộc cả hai vùng. Nó \textbf{không} tự động suy ra dynamic feasibility khi còn có control bounds, phân bổ thời gian, ràng buộc nonholonomic, hoặc động lực học phi tuyến. Vì vậy, một đường đi rời rạc từ $\sigma$ đến $\tau$ chỉ nên được hiểu là một \textbf{skeleton hình học khả dĩ}; tính khả thi động lực học sẽ được quyết định bởi các ràng buộc continuous-time ở các mục sau.

\section{Biến quyết định của mô hình}

Ta chia biến quyết định thành ba lớp.

\subsection{Biến rời rạc: chọn đường đi trên đồ thị}

Với mỗi cạnh $e=(u,v)\in E$, đặt
\begin{equation}
y_{uv}\in\{0,1\}.
\label{eq:ydef}
\end{equation}
Ý nghĩa:
\begin{itemize}
    \item $y_{uv}=1$: cạnh $(u,v)$ được chọn, nghĩa là chuỗi vùng đi từ $u$ sang $v$,
    \item $y_{uv}=0$: cạnh đó không được dùng.
\end{itemize}

Với mỗi đỉnh vùng $v\in V_r$, đặt thêm biến kích hoạt đỉnh
\begin{equation}
p_v\in\{0,1\}.
\label{eq:pdef}
\end{equation}
Ý nghĩa:
\begin{itemize}
    \item $p_v=1$: vùng $Q_v$ thực sự được quỹ đạo đi qua,
    \item $p_v=0$: vùng $Q_v$ không xuất hiện trong nghiệm.
\end{itemize}


Biến $p_v$ về nguyên tắc có thể được suy ra từ các biến cạnh $y_{uv}$ thông qua
\[
p_v=\sum_{(u,v)\in E} y_{uv}=\sum_{(v,w)\in E} y_{vw}, \qquad v\in V_r.
\]
Tuy nhiên, việc giữ riêng biến kích hoạt đỉnh $p_v$ giúp diễn đạt gọn hơn các ràng buộc và chi phí gắn trực tiếp với vùng $Q_v$ và với local trajectory segment tương ứng.


\subsection{Biến continuous cục bộ trên mỗi vùng}

Với mỗi vùng $v\in V_r$, ta gắn một \textbf{local trajectory segment} trên miền chuẩn hóa $\tau\in[0,1]$. Các biến continuous cục bộ là:
\begin{align}
& s_v^- \in \mathbb R^{n_x} && \text{: entry shooting state của vùng } v, \nonumber\\
& s_v^+ \in \mathbb R^{n_x} && \text{: exit state của vùng } v, \nonumber\\
& \Delta_v \ge 0 && \text{: thời lượng local của đoạn trong vùng } v, \nonumber\\
& w_v \in \mathbb R^{m_v} && \text{: tham số điều khiển trên vùng } v. \nonumber
\end{align}

Ngoài ra, từ các biến này ta sinh ra quỹ đạo cục bộ
\[
x_v:[0,1]\to\mathbb R^{n_x},
\qquad
u_v:[0,1]\to\mathbb R^{n_u},
\]
qua một local initial value problem sẽ được định nghĩa ở phần sau.

\paragraph{Giải thích.}
Trong direct multiple shooting, ta không quyết định trực tiếp toàn bộ trạng thái trên toàn chân trời thời gian chỉ từ đầu mút, mà thay vào đó mỗi đoạn thời gian có một \emph{ initial value} riêng $s_v^-$.

\subsection{Biến liên kết giữa hai vùng liên tiếp}

Với mỗi cạnh nội bộ $(u,v)\in E_{rr}$, đặt một biến interface
\begin{equation}
z_{uv}\in\mathbb R^{n_x}.
\label{eq:zdef}
\end{equation}
Ý nghĩa:
\begin{itemize}
    \item nếu cạnh $(u,v)$ được chọn, thì $z_{uv}$ là trạng thái vật lý tại kết nối giữa hai đoạn cục bộ trong vùng $u$ và vùng $v$,
    \item do đó $z_{uv}$ phải nằm trong phần giao $Q_u\cap Q_v$.
\end{itemize}

\section{Ràng buộc network flow cho phần rời rạc}

\subsection{Ràng buộc nguồn và đích}

Ta yêu cầu đúng một cạnh đi ra khỏi nguồn và đúng một cạnh đi vào đích:
\begin{equation}
\sum_{(\sigma,v)\in E} y_{\sigma v}=1,
\qquad
\sum_{(v,\tau)\in E} y_{v\tau}=1.
\label{eq:source_sink_flow}
\end{equation}

Ý nghĩa: nghiệm phải bắt đầu ở một vùng đầu và kết thúc ở một vùng cuối.

\subsection{Bảo toàn luồng tại mỗi vùng}

Với mỗi $v\in V_r$, ta áp
\begin{equation}
\sum_{(u,v)\in E} y_{uv}
=
\sum_{(v,w)\in E} y_{vw}
=
p_v.
\label{eq:flow_balance}
\end{equation}

Ý nghĩa của từng thành phần trong \eqref{eq:flow_balance}:
\begin{itemize}
    \item $\sum_{(u,v)} y_{uv}=p_v$: nếu vùng $v$ được dùng thì phải có đúng một cạnh đi vào,
    \item $\sum_{(v,w)} y_{vw}=p_v$: nếu vùng $v$ được dùng thì phải có đúng một cạnh đi ra,
    \item nếu $p_v=0$ thì không cạnh nào vào/ra vùng $v$ được kích hoạt.
\end{itemize}

\subsection{Ràng buộc loại chu trình}

\begin{equation}
\sum_{e\in I_v^{\mathrm{out}}} y_e \le 1,
\qquad \forall v\in V_r,
\label{eq:acyclicity}
\end{equation}
trong đó $I_v^{\mathrm{out}}$ là tập các cạnh đi ra từ đỉnh $v$.

\section{Mô hình động lực học cục bộ theo direct multiple shooting}

\subsection{Tham số hóa điều khiển trên từng vùng}

Với mỗi vùng $v\in V_r$, điều khiển local được tham số hóa bởi một vector hữu hạn chiều $w_v$ dưới dạng
\begin{equation}
u_v(\tau)=\mathcal U_v(\tau;w_v),
\qquad \tau\in[0,1].
\label{eq:control_param}
\end{equation}

Ví dụ, $\mathcal U_v(\tau;w_v)$ có thể là:
\begin{itemize}
    \item điều khiển hằng trên đoạn,
    \item điều khiển piecewise constant,
    \item spline hoặc polynomial bậc thấp.
\end{itemize}

Ví dụ, giả sử $u_v(\tau)\in\mathbb R^2$ là vận tốc phẳng của robot trong vùng $v$. Nếu chọn điều khiển piecewise constant với hai đoạn con, ta có thể viết
\[
u_v(\tau)=
\begin{cases}
u_{v,1}, & 0\le \tau < \tfrac12,\\
u_{v,2}, & \tfrac12 \le \tau \le 1,
\end{cases}
\qquad
w_v=(u_{v,1},u_{v,2}).
\]
Khi đó, thay vì tối ưu trên toàn bộ một hàm $u_v(\cdot)$, ta chỉ cần tối ưu trên hai vector điều khiển $u_{v,1}$ và $u_{v,2}$.

Ý tưởng của \eqref{eq:control_param} là thay bài toán tối ưu trên không gian hàm $u_v(\cdot)$, vốn là một bài toán vô hạn chiều, bằng bài toán tối ưu trên một vector hữu hạn chiều $w_v$. Nói cách khác, thay vì xem $u_v$ là một hàm tùy ý trên $[0,1]$, ta chỉ cho phép nó thuộc một họ hàm đã chọn trước và được mô tả hoàn toàn bởi các tham số trong $w_v$. Khi đó, một khi $w_v$ được cố định thì toàn bộ tín hiệu điều khiển $u_v(\tau)$ trên đoạn thuộc vùng $v$ cũng được xác định.

Việc tham số hóa này có hai vai trò chính:
\begin{itemize}
    \item Nó giúp đưa bài toán optimal control về dạng nonlinear program hữu hạn chiều, là dạng phù hợp với direct multiple shooting và với các solver số thực tế.
    \item Nó cho phép điều chỉnh mức độ chi tiết của lớp điều khiển: nếu chọn điều khiển hằng trên đoạn thì mô hình gọn và tốn ít chi phí tính toán hơn, nhưng kém linh hoạt hơn; nếu chọn piecewise constant với nhiều đoạn con, hoặc spline hay đa thức bậc thấp, thì quỹ đạo biểu diễn được sẽ phong phú hơn, đổi lại số biến quyết định và độ khó tính toán cũng tăng lên.
\end{itemize}

\subsection{Local IVP trên từng vùng}

Sau khi chuẩn hóa thời gian về $[0,1]$, local trajectory trong vùng $v$ được xác định bởi bài toán giá trị đầu
\begin{equation}
\begin{cases}
\dfrac{d x_v}{d\tau}(\tau)=\Delta_v\,f\bigl(x_v(\tau),u_v(\tau)\bigr), & \tau\in[0,1],\\[4pt]
x_v(0)=s_v^-,\\
u_v(\tau)=\mathcal U_v(\tau;w_v).
\end{cases}
\label{eq:local_ivp}
\end{equation}

Từ nghiệm của \eqref{eq:local_ivp}, ta định nghĩa endpoint map
\begin{equation}
F_v(s_v^-,w_v,\Delta_v):=x_v(1).
\label{eq:Fv}
\end{equation}

Ánh xạ $F_v$ được gọi là endpoint map, vì nó ánh xạ dữ liệu đầu của một đoạn động lực học gồm trạng thái đầu $s_v^-$, tham số điều khiển $w_v$ và thời lượng $\Delta_v$ thành trạng thái cuối thu được sau khi giải bài toán giá trị đầu \eqref{eq:local_ivp} trên toàn đoạn chuẩn hóa $[0,1]$.

Khi đó, ràng buộc defect của multiple shooting là
\begin{equation}
s_v^+ - F_v(s_v^-,w_v,\Delta_v)=0,
\qquad \forall v\in V_r.
\label{eq:defect}
\end{equation}

\paragraph{Ý nghĩa.}
Ràng buộc \eqref{eq:defect} nói rằng exit state của vùng $v$ phải đúng bằng endpoint của local IVP bắt đầu từ entry shooting state $s_v^-$. Đây là continuity condition trong direct multiple shooting, đúng với công thức chuẩn kiểu
\[
s_{i+1}=x_i(t_{i+1};s_i,q_i)
\]


\section{Ràng buộc an toàn hình học và vai trò của log-barrier}

Nếu vùng $v$ được kích hoạt, điều kiện mong muốn ở mức continuous time là local trajectory của vùng đó nằm trong vùng lồi tương ứng:
\begin{equation}
x_v(\tau)\in Q_v,
\qquad \forall \tau\in[0,1],
\quad \text{khi }p_v=1.
\label{eq:region_constraint}
\end{equation}
Vì $Q_v$ là vùng an toàn, nếu mỗi local segment nằm trong vùng được gán cho nó thì quỹ đạo toàn cục nằm trong hợp các safe regions đã chọn. Tuy nhiên, điều kiện \eqref{eq:region_constraint} là một ràng buộc continuous-time. Khi đưa vào solver hữu hạn chiều, ta không thể kiểm tra trực tiếp mọi $\tau\in[0,1]$ nếu không có một transcription đặc biệt cho ràng buộc liên tục.

Trong bản formulation cũ, điều kiện an toàn thường được xấp xỉ bằng cách kiểm tra tại một số mesh points hoặc ``red points''. Cách làm này có ý nghĩa như một kiểm tra rời rạc, nhưng không nên xem nó là cơ chế bảo đảm an toàn chính. Lý do là một tập hữu hạn điểm kiểm tra có thể bỏ sót vi phạm xảy ra giữa hai điểm liên tiếp. Trong ngôn ngữ multiple shooting, đây chính là hiện tượng local trajectory được sinh bởi tích phân động lực học có thể ``đi xuyên biên'' trong một shooting interval, sau đó quay lại vùng lồi trước khi chạm điểm kiểm tra tiếp theo.

Vì vậy, trong formulation mới, trọng tâm của safety được chuyển từ việc chỉ áp hard constraints tại các red points sang việc thêm một thành phần \emph{log-barrier} vào local running cost. Với mỗi vùng lồi được mô tả bởi
\begin{equation}
Q_v=\{x\in\mathbb R^{n_x}: A_vx\le b_v\},
\label{eq:polytope_region}
\end{equation}
log-barrier sẽ phạt rất mạnh khi trạng thái tiến sát một mặt biên của $Q_v$. Nhờ đó, optimizer có động lực giữ local trajectory nằm sâu hơn bên trong vùng lồi, thay vì chỉ vừa đủ thỏa mãn điều kiện tại một vài điểm rời rạc.

Các hard constraints tại một số điểm evaluation vẫn có thể được giữ lại, nhưng vai trò của chúng cần được hiểu đúng. Chúng không còn là toàn bộ cơ chế safety, mà chủ yếu là điều kiện miền xác định và điều kiện số học để log-barrier hợp lệ. Cụ thể, nếu $\tau_{v,q}$ là các quadrature/evaluation points được dùng để tính chi phí, ta có thể yêu cầu
\begin{equation}
(b_v)_j-(A_vx_v(\tau_{v,q}))_j\ge \delta_{\log}>0,
\qquad j=1,\dots,m_v,
\quad q\in\mathcal I_v^{\mathrm{eval}}.
\label{eq:barrier_domain_constraint}
\end{equation}
Tham số $\delta_{\log}$ là một ngưỡng dương nhỏ nhằm tránh $\log(0)$ hoặc log của số âm. Nếu vẫn dùng safety margin trong hiện thực, margin này nên được diễn giải là một hỗ trợ cho numerical stability và feasibility của barrier, không phải là trung tâm của phương pháp.

Nếu cần giữ ký hiệu $\mathcal S_v^{\mathrm{safe}}$, trong báo cáo này ta định nghĩa lại nó là tập các điều kiện miền xác định của log-barrier tại các điểm evaluation:
\begin{equation}
\mathcal S_v^{\mathrm{safe}}:=
\left\{
(b_v)_j-(A_vx_v(\tau_{v,q}))_j\ge \delta_{\log}
:\ j=1,\dots,m_v,\ q\in\mathcal I_v^{\mathrm{eval}}
\right\}.
\label{eq:safe_set_redefined}
\end{equation}
Cách viết này tránh hiểu nhầm rằng $\mathcal S_v^{\mathrm{safe}}\le 0$ là một sampling certificate đủ để bảo đảm continuous-time safety. Nó chỉ là phần kiểm tra hữu hạn điểm để đảm bảo các giá trị log-barrier được tính trong miền hợp lệ.

\section{Kiến thức nền tảng về log-barrier}

\subsection{Bất đẳng thức tuyến tính và slack}

Xét một polytope tổng quát
\begin{equation}
Q=\{x\in\mathbb R^n:Ax\le b\},
\label{eq:generic_polytope}
\end{equation}
trong đó $A\in\mathbb R^{m\times n}$ và $b\in\mathbb R^m$. Gọi $a_j^\top$ là hàng thứ $j$ của $A$, mỗi bất đẳng thức của polytope có dạng
\begin{equation}
a_j^\top x\le b_j,\qquad j=1,\dots,m.
\label{eq:single_halfspace}
\end{equation}
Ta định nghĩa slack của mặt biên thứ $j$ tại trạng thái $x$ là
\begin{equation}
s_j(x)=b_j-a_j^\top x.
\label{eq:slack_generic}
\end{equation}
Điểm $x$ nằm trong miền trong của $Q$ khi và chỉ khi tất cả các slack đều dương:
\begin{equation}
s_j(x)>0,
\qquad j=1,\dots,m.
\label{eq:interior_slack}
\end{equation}
Ngược lại, nếu một slack bằng $0$ thì $x$ nằm trên biên tương ứng; nếu một slack âm thì $x$ đã nằm ngoài polytope theo bất đẳng thức đó.

\subsection{Định nghĩa log-barrier}

Log-barrier chuẩn của polytope \eqref{eq:generic_polytope} được định nghĩa bởi
\begin{equation}
\phi(x)=-\sum_{j=1}^{m}\log\bigl(b_j-a_j^\top x\bigr)
= -\sum_{j=1}^{m}\log s_j(x).
\label{eq:generic_log_barrier}
\end{equation}
Ta chứng minh trực tiếp rằng $\phi(x)$ chỉ xác định trên miền trong của polytope. Thật vậy, theo định nghĩa
\begin{equation}
\phi(x)=-\sum_{j=1}^{m}\log\bigl(s_j(x)\bigr),
\end{equation}
và hàm $\log(\cdot)$ chỉ xác định khi đối số dương. Do đó, nếu tồn tại $j$ sao cho $s_j(x)\le 0$ thì $\phi(x)$ không xác định. Ngược lại, nếu $s_j(x)>0$ với mọi $j$ thì mọi hạng trong tổng đều hữu hạn, suy ra $\phi(x)$ xác định. Hơn nữa, với mỗi $j$ cố định, ta có $\log(s_j(x))\to -\infty$ khi $s_j(x)\to 0^+$, nên
\begin{equation}
-\log(s_j(x))\to +\infty.
\end{equation}
Vì vậy, khi $x$ tiến sát một mặt biên của polytope, ít nhất một slack tiến về $0^+$ và giá trị barrier tăng không bị chặn; còn khi $x$ nằm sâu trong miền trong thì các slack lớn hơn, nên từng hạng $-\log(s_j(x))$ nhỏ hơn và tổng $\phi(x)$ giảm.

Tính chất quan trọng này khiến log-barrier phù hợp để phạt trạng thái tiến sát biên. Thay vì chỉ báo vi phạm sau khi trạng thái đã vượt khỏi vùng, barrier tạo ra một chi phí tăng nhanh ngay khi trạng thái tiến gần biên từ phía trong.

\subsection{Ý nghĩa trong optimal control}

Với một optimal control problem có ràng buộc trạng thái
\begin{equation}
x(t)\in Q,
\qquad t\in[0,T],
\label{eq:ocp_state_in_Q}
\end{equation}
thay vì chỉ viết hard constraint $Ax(t)\le b$, ta có thể thêm barrier vào running cost:
\begin{equation}
\ell_{\mathrm{bar}}(x(t),u(t))
=
\ell(x(t),u(t))+\mu\,\phi(x(t)),
\qquad \mu>0.
\label{eq:bar_running_cost_generic}
\end{equation}
Hệ số $\mu$ điều chỉnh mức ảnh hưởng của barrier. Khi $\mu$ lớn, nghiệm bị đẩy xa biên mạnh hơn; khi $\mu$ nhỏ, nghiệm có thể tiến gần biên hơn để tối ưu các thành phần chi phí khác. Trong giai đoạn hiện tại, do chưa quan sát thấy nghiệm bị kéo về analytic center hoặc center của vùng lồi quá mạnh, formulation ưu tiên log-barrier chuẩn. Các biến thể như relaxed hoặc recentered log-barrier có thể được xem xét sau nếu xuất hiện vấn đề về infeasibility, độ lệch nghiệm quá lớn, hoặc cần một formulation trơn được mở rộng ra ngoài miền trong.

\subsection{Liên hệ với hiện tượng đi xuyên biên trong multiple shooting}

Trong direct multiple shooting, local trajectory trên mỗi vùng được sinh ra bằng cách tích phân động lực học từ shooting state $s_v^-$. Vì vậy, ngay cả khi entry state, exit state hoặc một vài mesh points nằm trong $Q_v$, phần trajectory giữa hai điểm kiểm tra vẫn có thể vượt ra ngoài vùng rồi quay lại. Đây là bản chất của inter-sample violation trong các phương pháp direct transcription.

Khi log-barrier được cộng vào running cost tại nhiều quadrature/evaluation points bên trong shooting interval, việc tiến sát biên trở nên đắt hơn về mặt objective. Cơ chế này không biến một tập hữu hạn điểm evaluation thành một chứng chỉ continuous-time tuyệt đối, nhưng nó giảm đáng kể xu hướng ``đi xuyên biên'' vì optimizer không chỉ phải thỏa điểm đầu/cuối mà còn phải trả chi phí lớn nếu trajectory đi gần biên tại các điểm bên trong interval. Do đó, phương pháp mới không phụ thuộc duy nhất vào việc tăng số lượng red points, dù việc tăng số điểm evaluation vẫn có thể cải thiện khả năng kiểm soát hành vi giữa các điểm rời rạc.

\section{Ràng buộc ghép nối giữa phần discrete và phần continuous}

Phần này buộc network-flow path và local multiple-shooting segments khớp nhau.

Với mỗi cạnh nội bộ $(u,v)\in E_{rr}$, ta áp logic sau:
\begin{equation}
y_{uv}=1 \Longrightarrow
\begin{cases}
z_{uv}\in Q_u\cap Q_v,\\
s_u^+=z_{uv},\\
s_v^-=z_{uv}.
\end{cases}
\label{eq:edge_coupling}
\end{equation}

Ý nghĩa từng dòng trong \eqref{eq:edge_coupling}:
\begin{itemize}
    \item $z_{uv}\in Q_u\cap Q_v$: điểm giao phải là một trạng thái hình học hợp lệ cho cả hai vùng,
    \item $s_u^+=z_{uv}$: đoạn local trong vùng $u$ kết thúc tại điểm nối,
    \item $s_v^-=z_{uv}$: đoạn local trong vùng $v$ bắt đầu từ đúng điểm nối ấy.
\end{itemize}

Như vậy, nếu cạnh $(u,v)$ được chọn thì hai local IVPs ở $u$ và $v$ được kết nối lại với nhau tại cùng một trạng thái vật lý.

Trong mô hình hóa cuối, perspective chỉ xử lý gọn phần membership $z_{uv}\in Q_u\cap Q_v$; còn hai đẳng thức $s_u^+=z_{uv}$ và $s_v^-=z_{uv}$ vẫn phải được bật/tắt rõ ràng bằng một cơ chế indicator tương thích với $y_{uv}$.

\section{Ràng buộc biên với source và target}

Ta cần gắn local trajectory đầu tiên với trạng thái đầu, và local trajectory cuối cùng với trạng thái đích.

Với mọi cạnh $(\sigma,v)\in E_s$:
\begin{equation}
y_{\sigma v}=1 \Longrightarrow s_v^-=x^{\mathrm s}.
\label{eq:source_bc}
\end{equation}

Với mọi cạnh $(v,\tau)\in E_t$:
\begin{equation}
y_{v\tau}=1 \Longrightarrow s_v^+=x^{\mathrm g}.
\label{eq:target_bc}
\end{equation}

Ý nghĩa: vùng đầu tiên trong chuỗi phải nhận đúng initial condition, còn vùng cuối cùng phải kết thúc ở goal state.


\section{Ép phương pháp log-barrier vào hiện thực hiện tại}

\subsection{Barrier của một vùng lồi}

Với mỗi vùng $v\in V_r$, giả sử mô tả polytope của vùng là
\begin{equation}
Q_v=\{x\in\mathbb R^{n_x}:A_vx\le b_v\},
\qquad A_v\in\mathbb R^{m_v\times n_x},\ b_v\in\mathbb R^{m_v}.
\label{eq:region_polytope_log}
\end{equation}
Gọi $a_{v,j}^\top$ là hàng thứ $j$ của $A_v$. Slack của mặt biên thứ $j$ tại trạng thái $x$ là
\begin{equation}
s_{v,j}(x)=b_{v,j}-a_{v,j}^\top x.
\label{eq:region_slack}
\end{equation}
Log-barrier của vùng $v$ được viết là
\begin{equation}
\phi_v(x)
=
-\sum_{j=1}^{m_v}\log\bigl((b_v)_j-(A_vx)_j\bigr)
=
-\sum_{j=1}^{m_v}\log\bigl(b_{v,j}-a_{v,j}^\top x\bigr).
\label{eq:region_log_barrier}
\end{equation}
Hàm $\phi_v$ chỉ được evaluate khi mọi slack trong \eqref{eq:region_slack} dương. Do đó, các điều kiện \eqref{eq:barrier_domain_constraint} là cần thiết trong transcription số.

\subsection{Running cost mới trong vùng}

Trong phần còn lại của báo cáo, running cost gốc được chọn theo dạng
\begin{equation}
\ell\bigl(x(t),u(t)\bigr)
:=
w_L\,\|\dot q(t)\|^2+w_E\,\|u(t)\|^2,
\qquad
q(t):=\pi\bigl(x(t)\bigr),
\qquad
w_L,w_E>0.
\label{eq:running_cost_choice}
\end{equation}
Ở đây, $q(t)$ là phần cấu hình hình học của trạng thái; nếu $x(t)$ bản thân đã là biến cấu hình thì có thể hiểu $q(t)=x(t)$. Thành phần $w_L\|\dot q(t)\|^2$ phạt độ lớn vận tốc dịch chuyển, còn $w_E\|u(t)\|^2$ phạt độ lớn điều khiển.

Khi đưa log-barrier vào vùng $v$, running cost local trở thành
\begin{equation}
\ell_v^{\mathrm{bar}}(x,u)
=
\ell(x,u)+\mu_v\phi_v(x)
=
w_L\|\dot q\|^2+w_E\|u\|^2+
\mu_v\phi_v(x),
\label{eq:bar_running_cost_region}
\end{equation}
trong đó $\mu_v>0$ là trọng số barrier của vùng $v$. Thành phần phạt thời gian $a\Delta_v$ không bị thay thế bởi log-barrier; nó vẫn là một hạng riêng trong local cost để tạo áp lực giảm tổng thời gian.

\subsection{Local cost mới trên miền chuẩn hóa}

Với local trajectory $x_v(\tau)$, local control $u_v(\tau)$, thời lượng $\Delta_v$ và $q_v(\tau):=\pi(x_v(\tau))$, local cost mới được định nghĩa là
\begin{equation}
\begin{aligned}
J_v^{\mathrm{bar}}(s_v^-,w_v,\Delta_v)
:={}& a\,\Delta_v \\
&+\int_0^1 \Delta_v\left(
w_L\left\|\frac{1}{\Delta_v}\frac{dq_v}{d\tau}(\tau)\right\|^2
+w_E\|u_v(\tau)\|^2
+\mu_v\phi_v\bigl(x_v(\tau)\bigr)
\right)d\tau .
\end{aligned}
\label{eq:local_cost_barrier}
\end{equation}
Hệ số $\Delta_v$ xuất hiện vì đoạn thời gian thật của vùng $v$ được chuẩn hóa về $\tau\in[0,1]$. Nếu $t=t_v^-+\Delta_v\tau$, thì $dt=\Delta_v d\tau$. Do đó, mọi running cost tích lũy theo thời gian thật đều phải được nhân với Jacobian $\Delta_v$ khi viết trên miền chuẩn hóa.

Cần lưu ý rằng $\dot q(t)$ trong thời gian thật tương ứng với
\begin{equation}
\dot q(t)=\frac{1}{\Delta_v}\frac{dq_v}{d\tau}(\tau).
\label{eq:qdot_scaled}
\end{equation}
Vì vậy, hạng vận tốc trong \eqref{eq:local_cost_barrier} có dạng $\left\|(1/\Delta_v)dq_v/d\tau\right\|^2$, còn toàn bộ running cost vẫn được nhân với $\Delta_v$ do phép đổi biến thời gian.

\subsection{Rời rạc hóa để hiện thực}

Trong hiện thực số, tích phân trong \eqref{eq:local_cost_barrier} được xấp xỉ bằng quadrature hoặc bằng các điểm evaluation sinh ra trong quá trình tích phân RK4. Giả sử trong vùng $v$ có các điểm evaluation
\begin{equation}
\tau_{v,q},\qquad q\in\mathcal I_v^{\mathrm{eval}},
\label{eq:evaluation_points}
\end{equation}
với trọng số quadrature $\omega_q$. Khi đó
\begin{equation}
\begin{aligned}
\widehat J_v^{\mathrm{bar}}
:={}& a\Delta_v \\
&+\sum_{q\in\mathcal I_v^{\mathrm{eval}}}\omega_q\Delta_v
\left[
w_L\left\|\frac{1}{\Delta_v}\frac{dq_v}{d\tau}(\tau_{v,q})\right\|^2
+w_E\|u_v(\tau_{v,q})\|^2
+\mu_v\phi_v\bigl(x_v(\tau_{v,q})\bigr)
\right].
\end{aligned}
\label{eq:discrete_barrier_cost}
\end{equation}
Barrier nên được evaluate tại nhiều điểm bên trong shooting interval, không chỉ tại entry state hoặc exit state. Nếu implementation đã dùng RK4, các intermediate stages của RK4 là các ứng viên tự nhiên để evaluate thêm barrier, vì chúng phản ánh trạng thái trung gian của local integration. Tăng số evaluation points giúp kiểm soát tốt hơn hành vi giữa hai điểm kiểm tra, nhưng đồng thời làm tăng chi phí tính toán và số hạng phi tuyến trong objective.

Để \eqref{eq:discrete_barrier_cost} hợp lệ, với mọi điểm evaluation cần áp điều kiện miền xác định
\begin{equation}
(b_v)_j-(A_vx_v(\tau_{v,q}))_j\ge \delta_{\log}>0,
\qquad
j=1,\dots,m_v,
\quad q\in\mathcal I_v^{\mathrm{eval}}.
\label{eq:discrete_domain_barrier_again}
\end{equation}
Đây là domain/numerical constraint cho log-barrier. Nó không nên được diễn giải là một chứng chỉ an toàn continuous-time hoàn chỉnh.

\subsection{Lựa chọn trọng số $\mu_v$}

Trọng số $\mu_v$ quyết định mức độ optimizer bị phạt khi local trajectory tiến gần biên vùng lồi. Một cách thực nghiệm hợp lý là bắt đầu với $\mu_v$ tương đối lớn để đẩy nghiệm vào phía trong vùng lồi, sau đó giảm dần nếu muốn cho phép nghiệm tiến gần biên hơn nhằm giảm thời gian hoặc năng lượng. Các mức có thể thử trong giai đoạn đầu là $10^{-2}$, $10^{-3}$, $10^{-4}$ và $10^{-5}$, tùy theo scale của các hạng chi phí còn lại.

Trong formulation hiện tại, ta ưu tiên log-barrier chuẩn vì mục tiêu chính là giảm hiện tượng đi xuyên biên giữa các shooting interval. Nếu về sau nghiệm bị kéo vào quá sâu trong vùng lồi, hoặc bài toán thường xuyên infeasible do yêu cầu slack dương quá nghiêm ngặt, khi đó có thể chuyển sang relaxed log-barrier hoặc recentered log-barrier.

\section{Hàm mục tiêu của mô hình sau khi thêm log-barrier}

\subsection{Objective dạng ý niệm}

Sau khi thay local cost $J_v$ bằng $J_v^{\mathrm{bar}}$, objective toàn cục ở mức ý niệm là
\begin{equation}
\min \sum_{v\in V_r}p_v\,J_v^{\mathrm{bar}}(s_v^-,w_v,\Delta_v).
\label{eq:global_cost_barrier_conceptual}
\end{equation}
Chỉ những vùng active mới đóng góp chi phí. Vì barrier cost nằm bên trong $J_v^{\mathrm{bar}}$, nó cũng chỉ áp dụng cho local trajectory của vùng được kích hoạt. Khi $p_v=0$, local block của vùng $v$ bị tắt theo logic on--off đã nêu ở phần perspective và không đóng góp vào objective.

\subsection{Objective dạng epigraph}

Dạng khuyến nghị cho solver là
\begin{equation}
\min \sum_{v\in V_r}\eta_v,
\label{eq:global_cost_barrier_epigraph}
\end{equation}
trong đó $\eta_v$ là biến epigraph (tức một upper bound) cho chi phí local đã bao gồm log-barrier. Logic bật/tắt được viết là
\begin{equation}
p_v=1\Longrightarrow
\eta_v\ge \widehat J_v^{\mathrm{bar}}(s_v^-,w_v,\Delta_v),
\label{eq:rho_active_barrier}
\end{equation}
\begin{equation}
p_v=0\Longrightarrow \eta_v=0.
\label{eq:rho_inactive_barrier}
\end{equation}

Giải thích ngắn: bất đẳng thức $\eta_v\ge\widehat J_v^{\mathrm{bar}}$ đặt $\eta_v$ làm một upper bound (epigraph) của chi phí local; vì objective tối thiểu hoá tổng $\eta_v$, tại nghiệm tối ưu sẽ có $\eta_v=\widehat J_v^{\mathrm{bar}}$ khi vùng được kích hoạt. Việc dùng biến epigraph tách rời chi phí khỏi phần dynamics/on--off giúp trình bày và triển khai với solver dạng MINLP rõ ràng hơn.

Với cách viết này, $\eta_v$ chứa chi phí thời gian, running cost động lực học, năng lượng điều khiển và log-barrier cost của vùng lồi. Nếu vùng không active, $\eta_v$ bị triệt tiêu và toàn bộ local block được neo về trạng thái off.

\section{Formulation thô ban đầu của bài toán ghép}

Gom các thành phần ý niệm lại, ta thu được formulation thô ban đầu sau:
\begin{align}
\min_{\substack{y,p,z,\\ \{s_v^-,s_v^+,w_v,\Delta_v\}_{v\in V_r}}}
\quad & \sum_{v\in V_r} p_v\,J_v^{\mathrm{bar}}(s_v^-,w_v,\Delta_v)
\label{eq:joint_problem_obj}\\
\text{s.t.}\quad
& \sum_{(\sigma,v)\in E} y_{\sigma v}=1,
\qquad
\sum_{(v,\tau)\in E} y_{v\tau}=1,
\label{eq:joint_problem_source_sink}\\
& \sum_{(u,v)\in E} y_{uv}
=
\sum_{(v,w)\in E} y_{vw}
=
p_v,
\qquad \forall v\in V_r,
\label{eq:joint_problem_flow}\\
& \sum_{e\in I_v^{\mathrm{out}}} y_e \le 1,
\qquad \forall v\in V_r,
\label{eq:joint_problem_acyc}\\
& x_v'(\tau)=\Delta_v f\bigl(x_v(\tau),u_v(\tau)\bigr),
\qquad \tau\in[0,1],
\quad \forall v\in V_r,
\label{eq:joint_problem_dyn}\\
& x_v(0)=s_v^-,
\qquad
u_v(\tau)=\mathcal U_v(\tau;w_v),
\quad \forall v\in V_r,
\label{eq:joint_problem_ivp}\\
& s_v^+ - F_v(s_v^-,w_v,\Delta_v)=0,
\qquad \forall v\in V_r,
\label{eq:joint_problem_defect}\\
& \mathcal S_v^{\mathrm{safe}} \text{ trong } \eqref{eq:safe_set_redefined} \text{ được thỏa},
\qquad \forall v\in V_r \text{ với } p_v=1,
\label{eq:joint_problem_region}\\
& y_{uv}=1 \Longrightarrow
\begin{cases}
z_{uv}\in Q_u\cap Q_v,\\
s_u^+=z_{uv},\\
s_v^-=z_{uv},
\end{cases}
\qquad \forall (u,v)\in E_{rr},
\label{eq:joint_problem_edge}\\
& y_{\sigma v}=1 \Longrightarrow s_v^-=x^{\mathrm s},
\qquad \forall (\sigma,v)\in E_s,
\label{eq:joint_problem_source}\\
& y_{v\tau}=1 \Longrightarrow s_v^+=x^{\mathrm g},
\qquad \forall (v,\tau)\in E_t,
\label{eq:joint_problem_target}\\
& y_{uv}\in\{0,1\},\qquad p_v\in\{0,1\},
\qquad \forall (u,v)\in E,\ \forall v\in V_r.
\label{eq:joint_problem_binary}
\end{align}


Trong \eqref{eq:joint_problem_obj}--\eqref{eq:joint_problem_binary}, ràng buộc loại chu trình được giữ lại để giới hạn số cạnh đi ra từ mỗi vùng, qua đó tránh việc lời giải rời rạc tạo vòng lặp không cần thiết.

\section{Perspective cho hình học lồi và động lực học}

\subsection{Vấn đề của formulation}
Nếu dùng trực tiếp bài toán \eqref{eq:joint_problem_obj}--\eqref{eq:joint_problem_binary}, mô hình thường gặp ba vướng mắc chính:
\begin{enumerate}
    \item objective kiểu
    \[
    \sum_{v\in V_r} p_v\,J_v^{\mathrm{bar}}(s_v^-,w_v,\Delta_v)
    \]
    nhìn trực quan như một cơ chế ``bật/tắt'' local cost, nhưng thực chất đang nhân biến nhị phân với một hàm local nói chung là phi lồi;
    \item các biến continuous ở những vùng không được chọn có thể vẫn ``trôi tự do'' nếu không có cơ chế deactivation minh bạch;
    \item nếu mô hình hóa ghép nối bằng tích nhị phân--liên tục (ví dụ $y_{uv}z_{uv}$), ta sẽ đưa thêm các bilinear constraints không cần thiết.
\end{enumerate}

Vì vậy, cách viết ổn hơn là tách mô hình thành hai lớp:
\begin{enumerate}
    \item \textbf{lớp hình học lồi/on--off}: dùng \emph{homogenization/perspective} cho các tập lồi gắn với đỉnh và cạnh;
    \item \textbf{lớp động lực học multiple shooting}: giữ riêng dưới dạng các ràng buộc phi tuyến được bật/tắt bằng indicator/disjunction, không cố ép chúng vào một perspective form khi bản thân chúng không lồi.
\end{enumerate}

\subsection{Perspective cho các tập lồi}
Giả sử $C\subset\mathbb R^n$ là một tập lồi đóng. Ta định nghĩa homogenization của $C$ bởi
\begin{equation}
\widetilde C:=\{(x,\lambda):\lambda>0,\ x/\lambda\in C\},
\qquad
\operatorname{cl}(\widetilde C)
\text{ là closure của nó.}
\label{eq:homogenized-set}
\end{equation}

Ý nghĩa của \eqref{eq:homogenized-set} là:
\begin{itemize}
    \item nếu $\lambda=1$ thì ta quay về đúng $x\in C$,
    \item nếu $\lambda=0$ thì ta thu được trạng thái ``off'' theo đúng closure của homogenization,
\end{itemize}

Áp dụng vào bài toán hiện tại, với mỗi vùng $v\in V_r$, ta có thể viết các ràng buộc hình học on--off dưới dạng
\begin{equation}
(s_v^-,p_v)\in \operatorname{cl}(\widetilde{Q}_v),
\qquad
(s_v^+,p_v)\in \operatorname{cl}(\widetilde{Q}_v).
\label{eq:perspective-vertex-states}
\end{equation}

Tương tự, với mỗi cạnh nội bộ $(u,v)\in E_{rr}$, thay vì chỉ viết
\[
z_{uv}\in Q_u\cap Q_v
\quad\text{khi } y_{uv}=1,
\]
ta viết trực tiếp
\begin{equation}
(z_{uv},y_{uv})
\in
\operatorname{cl}\!\bigl(\widetilde{(Q_u\cap Q_v)}\bigr).
\label{eq:perspective-edge-interface}
\end{equation}
Khi đó:
\begin{itemize}
    \item nếu $y_{uv}=1$ thì \eqref{eq:perspective-edge-interface} khôi phục đúng điều kiện $z_{uv}\in Q_u\cap Q_v$;
    \item nếu $y_{uv}=0$ thì biến $z_{uv}$ bị ``triệt tiêu'' theo closure của homogenization và không còn là một biến tự do có ý nghĩa vật lý độc lập.
\end{itemize}

Perspective ở đây chỉ xử lý gọn phần \emph{membership} trên các tập lồi. Nó không tự động mã hóa các đẳng thức ghép nối $s_u^+=z_{uv}$ và $s_v^-=z_{uv}$; hai đẳng thức này vẫn phải được viết tường minh dưới dạng on--off constraints gắn với $y_{uv}$.

\subsection{Bật/tắt local multiple-shooting block}
Với mỗi vùng $v$, local multiple-shooting block vẫn phải giữ nguyên dưới dạng
\begin{equation}
\begin{cases}
\dfrac{d x_v}{d\tau}(\tau)=\Delta_v\,f\bigl(x_v(\tau),u_v(\tau)\bigr), & \tau\in[0,1],\\[4pt]
x_v(0)=s_v^-,\\
u_v(\tau)=\mathcal U_v(\tau;w_v),\\[2pt]
s_v^+ - F_v(s_v^-,w_v,\Delta_v)=0.
\end{cases}
\label{eq:ms-block-repeat}
\end{equation}
Khối \eqref{eq:ms-block-repeat} nói chung là \textbf{phi lồi}, vì $F_v$ là endpoint map sinh ra bởi tích phân một ODE và phụ thuộc phi tuyến vào các biến quyết định. Paper \emph{Motion Planning around Obstacles with Convex Optimization} cũng nhấn mạnh rằng nếu thêm continuous-time dynamics vào mô hình thì các equality constraints tương ứng trở nên nonconvex, ngay cả với linear control systems khi shape và timing cùng là biến quyết định \cite{marcucci2022}.

Do đó, perspective \textbf{không} thể thay thế phần này. Cách làm đúng là giữ nó như một \textbf{nonlinear dynamic block} và bật/tắt bằng indicator/disjunction.


Với mỗi vùng $v\in V_r$, ta dùng logic sau:
\begin{equation}
p_v=1
\Longrightarrow
\begin{cases}
\Delta_v \ge \underline \Delta_v >0,\\
\text{local IVP } \eqref{eq:local_ivp} \text{ được kích hoạt},\\
s_v^+ - F_v(s_v^-,w_v,\Delta_v)=0,\\
\mathcal S_v^{\mathrm{safe}} \text{ trong } \eqref{eq:safe_set_redefined} \text{ được thỏa},\\
\eta_v \ge \widehat J_v^{\mathrm{bar}}(s_v^-,w_v,\Delta_v),
\end{cases}
\label{eq:active-vertex-block}
\end{equation}
trong đó $\eta_v$ là biến epigraph cho local cost trên vùng $v$, $\widehat J_v^{\mathrm{bar}}$ là một xấp xỉ quadrature của chi phí đã bao gồm log-barrier, còn $\mathcal S_v^{\mathrm{safe}}$ được hiểu là các điều kiện miền xác định của log-barrier tại các điểm evaluation/quadrature, như trong \eqref{eq:safe_set_redefined}.

Ngược lại, khi vùng $v$ không được dùng, ta triệt tiêu toàn bộ block:
\begin{equation}
p_v=0
\Longrightarrow
\begin{cases}
\Delta_v = 0,\\
\eta_v = 0,\\
s_v^- = 0,\\
s_v^+ = 0,\\
w_v = \bar w_v,
\end{cases}
\label{eq:inactive-vertex-block}
\end{equation}
trong đó $s_v^\pm=0$ là lựa chọn off-state tương thích với \eqref{eq:perspective-vertex-states}, còn $\bar w_v$ là một giá trị neo cố định đã chọn trước, thường lấy bằng $0$.

Ý nghĩa của \eqref{eq:active-vertex-block}--\eqref{eq:inactive-vertex-block} là:
\begin{itemize}
    \item nếu vùng active thì multiple-shooting dynamics và local cost thật sự được áp lên;
    \item nếu vùng inactive thì mọi biến local bị khóa vào trạng thái off và không thể ``trôi tự do'' trong mô hình.
\end{itemize}

\subsection{Mô hình cuối}
Không nên xem
\[
\sum_v p_v J_v^{\mathrm{bar}}(s_v^-,w_v,\Delta_v)
\]
là formulation solver-ready cuối cùng. Sau khi chỉnh các điểm ở trên, mô hình nên được trình bày theo bốn lớp rõ ràng:
\begin{enumerate}
    \item \textbf{lớp discrete path}:
    \begin{align*}
    & \eqref{eq:source_sink_flow},\ \eqref{eq:flow_balance},\ \eqref{eq:acyclicity};
    \end{align*}
    \item \textbf{lớp convex/on--off geometry}:
    \begin{align*}
    & (s_v^-,p_v),(s_v^+,p_v) \in \operatorname{cl}(\widetilde{Q}_v),\qquad \forall v\in V_r,\\
    & (z_{uv},y_{uv})\in \operatorname{cl}\!\bigl(\widetilde{(Q_u\cap Q_v)}\bigr),\qquad \forall (u,v)\in E_{rr};
    \end{align*}
    \item \textbf{lớp ghép nối on--off}:
    \begin{align*}
    & y_{uv}=1 \Longrightarrow s_u^+=z_{uv},\qquad y_{uv}=1 \Longrightarrow s_v^-=z_{uv},
    \qquad \forall (u,v)\in E_{rr},\\
    & y_{\sigma v}=1 \Longrightarrow s_v^-=x^{\mathrm s},\qquad \forall (\sigma,v)\in E_s,\\
    & y_{v\tau}=1 \Longrightarrow s_v^+=x^{\mathrm g},\qquad \forall (v,\tau)\in E_t;
    \end{align*}
    \item \textbf{lớp nonlinear multiple-shooting dynamics}:
    \begin{align*}
    & p_v=1 \Longrightarrow
    \begin{cases}
    \text{local IVP } \eqref{eq:local_ivp},\\
    s_v^+ - F_v(s_v^-,w_v,\Delta_v)=0,\\
    \mathcal S_v^{\mathrm{safe}} \text{ trong } \eqref{eq:safe_set_redefined} \text{ được thỏa},\\
    \eta_v \ge \widehat J_v^{\mathrm{bar}}(s_v^-,w_v,\Delta_v),
    \end{cases}
    \\
    & p_v=0 \Longrightarrow
    \begin{cases}
    \Delta_v=0,\ \eta_v=0,\\
    s_v^- = s_v^+ = 0,\\
    w_v=\bar w_v.
    \end{cases}
    \end{align*}
\end{enumerate}
Khi đó objective toàn cục được viết gọn hơn thành
\begin{equation}
\min \sum_{v\in V_r} \eta_v,
\label{eq:recommended-clean-objective}
\end{equation}
trong đó:
\begin{itemize}
    \item $\eta_v$ chứa chi phí thời gian, local dynamic cost, energy cost và log-barrier cost của multiple-shooting segment trên vùng $v$;
    \item block chi phí này được bật/tắt bằng indicator/disjunction đồng bộ với activation variable $p_v$.
\end{itemize}

\paragraph{Kết luận}
Perspective function được dùng để xử lý phần on--off convex của mô hình, cụ thể là:
\begin{itemize}
    \item triệt tiêu các biến continuous tại các vùng/cạnh không active,
    \item viết gọn các interface constraints trên các tập lồi giao nhau.
\end{itemize}
Ở đây, ``interface constraints'' là các ràng buộc ghép nối tại điểm chuyển tiếp giữa hai vùng liên tiếp, chẳng hạn điều kiện $z_{uv}\in Q_u\cap Q_v$ cùng với $s_u^+=z_{uv}$ và $s_v^-=z_{uv}$. Dùng perspective cho biến kết nối giúp biểu diễn gọn logic ``cạnh active thì điểm nối phải nằm trong phần giao, cạnh inactive thì biến kết nối bị triệt tiêu'', thay vì phải viết nhiều ràng buộc on--off rời rạc hoặc sinh thêm các tích nhị phân--liên tục không cần thiết.
Tuy nhiên, perspective \textbf{không} loại bỏ được bản chất phi lồi của multiple shooting, cũng không thay thế cho cycle-elimination, defect constraints hay các điều kiện miền xác định của log-barrier. Nói cách khác, sau khi chỉnh mô hình đúng cách, ta vẫn đang làm việc với một \textbf{mixed-integer nonlinear optimal-control framework} ở mức continuous time.

\section{Phân loại bài toán}

Mô hình trên \textbf{không còn là MISOCP} như GCS-B\'ezier. Lý do là các ánh xạ
\[
F_v(s_v^-,w_v,\Delta_v)
\]
được tạo ra bởi việc tích phân một ODE phi tuyến nói chung, và các ràng buộc vùng liên tục theo thời gian cũng không còn là thuần convex algebraic constraints.

Do đó, cách gọi chính xác hơn là:
\begin{equation*}
\boxed{
\begin{array}{c}
\text{Continuous-time level: Mixed-Integer Nonlinear Optimal Control Problem (MIOCP),}\\
\text{after finite-dimensional transcription: Mixed-Integer Nonlinear Program (MINLP).}
\end{array}}
\end{equation*}

\begin{thebibliography}{99}

\bibitem{bockplitt}
H.~G. Bock and K.~J. Plitt,
\newblock \emph{A Multiple Shooting Algorithm for Direct Solution of Optimal Control Problems},
\newblock Proceedings of the 9th IFAC World Congress, Budapest, 1984.

\bibitem{diehl2005}
M. Diehl, H.~G. Bock, H. Diedam, and P.-B. Wieber,
\newblock \emph{Fast Direct Multiple Shooting Algorithms for Optimal Robot Control},
\newblock in \emph{Fast Motions in Biomechanics and Robotics}, Springer, 2006.

\bibitem{marcucci2021}
T. Marcucci, J. Umenberger, P.~A. Parrilo, and R. Tedrake,
\newblock \emph{Shortest Paths in Graphs of Convex Sets},
\newblock SIAM Journal on Optimization, 2023.

\bibitem{marcucci2022}
T. Marcucci, M. Petersen, D. von Wrangel, and R. Tedrake,
\newblock \emph{Motion Planning around Obstacles with Convex Optimization},
\newblock The International Journal of Robotics Research, 2023.

\bibitem{marcucciThesis}
T. Marcucci,
\newblock \emph{Graphs of Convex Sets with Applications to Optimal Control and Motion Planning},
\newblock PhD thesis, Massachusetts Institute of Technology, 2024.


\bibitem{schiela2008}
A. Schiela,
\newblock \emph{Barrier Methods for Optimal Control Problems with State Constraints},
\newblock ZIB-Report 07-07, 2008.

\bibitem{feller2015}
C. Feller and C. Ebenbauer,
\newblock \emph{Relaxed Logarithmic Barrier Function Based Model Predictive Control of Linear Systems},
\newblock arXiv:1503.03314, 2015.

\bibitem{elango2025}
P. Elango, D. Luo, A. G. Kamath, S. Uzun, T. Kim, and B. Açıkmeşe,
\newblock \emph{Continuous-time successive convexification for constrained trajectory optimization},
\newblock Automatica, 180:112464, 2025.

\end{thebibliography}

\end{document}
